\documentclass[10pt,a4paper]{amsart}
\usepackage[margin=19mm]{geometry}
\usepackage{amsmath,amssymb,mathtools}
\usepackage[scr=rsfs,cal=euler]{mathalfa}
\usepackage{xfrac}
\usepackage[unicode,hidelinks,bookmarks=false]{hyperref}
\newcommand{\ie}{i.e.\ }
\newcommand{\cf}{cf.\ }
\newcommand{\resp}{respectively\ }
\newcommand{\Subsection}{\subsection}
\providecommand{\nobreakdash}{\nobreak\hbox{-}}
\setlength{\emergencystretch}{2em}
\multlinegap0pt
\usepackage{bm}
\usepackage{amssymb}
\usepackage{xcolor}
\newtheorem{assumption}{Assumption}
\newtheorem{proposition}{Proposition}
\newtheorem{lemma}{Lemma}
\usepackage{cleveref}
\crefname{assumption}{Assumption}{Assumptions}
\crefname{proposition}{Proposition}{Propositions}
\crefname{lemma}{Lemma}{Lemmas}
\newcommand{\bfitgreek}[1]{\bm{\mathit{#1}}}
\title{Stabilization-Free \(\mathbf{H}\left(\operatorname{\mathbf{curl}}\right)\) and \(\mathbf{H}\left(\operatorname{div}\right)\)-Conforming Virtual Element Method}
\author{Yuxuan Liao, Xue Feng, Yidong Huang}
\date{Source: arXiv:2501.15168v2 (CC BY 4.0)}
\begin{document}
\maketitle

\noindent\emph{Source and licence.} Stabilization-Free \(\mathbf{H}\left(\operatorname{\mathbf{curl}}\right)\) and \(\mathbf{H}\left(\operatorname{div}\right)\)-Conforming Virtual Element Method. Version arXiv:2501.15168v2 (CC BY 4.0), distributed by arXiv under the Creative Commons Attribution 4.0 International licence, http://creativecommons.org/licenses/by/4.0/, as recorded in the arXiv licence metadata for this exact version and shown on the abstract page https://arxiv.org/abs/2501.15168v2. Full licence notice: \texttt{licenses/arxiv-2501.15168-CC-BY-4.0.txt}.\par\smallskip

\noindent\textit{Editorial scope.} Counted unit: the article's lemma lem:sfsubbound (source0.tex lines 588--596). The article's complete original proof of it is delivered with it (source0.tex lines 598--632, one \texttt{proof} environment of 35 lines). No mathematics is written, repaired or re-derived: the statement and the proof are the article's own lines, in the article's own notation, and no retained line is substituted or re-flowed.  Retained uncounted as the source-local context the proof needs: lines 281--296; lines 430--432.  Nothing was omitted from any retained span, so the package carries no omission record.  No retained line contains a citation, so the wrapper carries no bibliography and no bibliography entry.  No retained line contains a figure environment, an image-inclusion command or drawing code, so no image is delivered and the package declares no asset dependency.  Editorial formatting only, in addition: an AMS \texttt{amsart} preamble that restates the article's own declarations and macros used by the retained lines, namely the environments \texttt{assumption}, \texttt{proposition}, \texttt{lemma} on separate counters, together with the standard packages those lines need.  The retained environments print in this document as Assumption 1, Proposition 1, Lemma 1 in order of appearance, whereas the article prints the same items with the numbering of its own full document.  The complete native source of the article as distributed in the arXiv e-print for this exact version is bundled unchanged as \texttt{sources/172157/source0.tex}.
\begin{assumption}[Mesh Regularity]
  \label{ass:mesh}There exists a uniform constant \(\rho>0\) such that, for every
  polyhedron \(K\in \mathcal{T}_{h}\),
  \begin{enumerate}
    \def\labelenumi{(\roman{enumi})}
    \item
          \(K\) is star-shaped with respect to a ball of radius
          \(\geq \rho h_K\);
    \item
          every face \(F\) of \(\partial K\) is star-shaped with respect to a
          disk of radius \(\geq \rho h_F\);
    \item
          for every face \(F\) of \(\partial K\), every edge \(E\) of
          \(\partial F\) satisfies \( h_E \geq \rho h_F \geq \rho^2 h_K\).
  \end{enumerate}
\end{assumption}

\begin{proposition}
  \label{prop:sf}\(\bfitgreek{\Pi}^{\mathrm{Sf}}_{k}\) is well-defined.
\end{proposition}

\begin{lemma}
  \label{lem:sfsubbound}The following bound holds true: \(\forall \boldsymbol{q}_{k+1}\in\mathbb{P}_{k+1}^{3}(K)\)
  \begin{equation}\label{eq:contrasf}
    \begin{aligned}
       & h^{\frac{1}{2}}_{K}\left\|\nabla \times \boldsymbol{q}_{k+1}\cdot \boldsymbol{n}\right\|_{\partial K}+\sup_{\boldsymbol{p}_{\mu_{\mathrm{r}}}\in \mathbb{P}_{\mu_{\mathrm{r}}}^{3}\left(K\right)}\frac{\left(\nabla \times \boldsymbol{q}_{k+1},\boldsymbol{x}_K \times \boldsymbol{p}_{\mu_{\mathrm{r}}}\right)_{K}}{\left\|\boldsymbol{x}_K \times \boldsymbol{p}_{\mu_{\mathrm{r}}}\right\|_{K}} \\
       & \gtrsim\left\|\nabla \times\boldsymbol{q}_{k+1}\right\|_{K}.
    \end{aligned}
  \end{equation}
\end{lemma}

\begin{proof}
  It suffices to prove the claim fixing \(h_{K}=1\), and then use a scaling argument. By the mesh assumption, there are concentric balls \(\tilde{B}\) and \(B\), with radius \(\rho\) and \(1+\rho\), satisfying:
  \begin{equation}
    \tilde{B}\subseteq K\subseteq B.
  \end{equation}

  Let \(S_{N_{\mathrm{v}}}\) denote the set of polytopes with \(N_{\mathrm{v}}\) vertices satisfying \cref{ass:mesh}. Setting the center of \(\tilde{B}\) and \(B\) at the origin, we can define an injective mapping \(\mathscr{M}:S_{N_{\mathrm{v}}}\to \mathbb{R}^{3N_{\mathrm{v}}}\) by:
  \begin{equation}
    K\mapsto\left[v_{1,x},v_{1,y},v_{1,z},\dots,v_{N_{\mathrm{v}},x},v_{N_{\mathrm{v}},y},v_{N_{\mathrm{v}},z}\right]^{\top},
  \end{equation}
  where \(v_{n,x},v_{n,y},v_{n,z}\) are the \(x,y,z\) coordinates of the \(n\)-th vertex. Note that we here assume the admissibility of the polytopes and well-defined order of the vertices, though it is not the case when \(S\) is only categorized by vertex number. But since the following proof is applicable to any further categorization of \(S_{N_{\mathrm{v}}}\), we resort to avoiding the tedious task.

  Now we will prove the claim by contradiction. If \cref{eq:contrasf} were false, there exists a sequence of polytopes \(\{K_{n}\in S_{N_{\mathrm{v}}}\}_{n\in \mathbb{N}}\), and a sequence of polynomials \(\{\nabla \times \boldsymbol{q}_{n}\in \nabla \times \mathbb{P}^{3}_{k+1}(B)\}_{n\in \mathbb{N}}\), such that:
  \begin{equation}\label{eq:pfcontrasf}
    \|\nabla \times \boldsymbol{q}_{n}\|_{K_{n}}=1,
    \|\nabla \times \boldsymbol{q}_{n}\cdot \boldsymbol{n}\|_{\partial K_{n}}\le\frac{1}{n},
    \sup_{\boldsymbol{p}_{\mu_{\mathrm{r}}}\in\mathbb{P}^{3}_{\mu_{\mathrm{r}}}\left(K_n\right)}
    \frac{(\nabla \times \boldsymbol{q}_{n},\,\boldsymbol{x}_K\times\boldsymbol{p}_{\mu_{\mathrm{r}}})_{K_{n}}}
    {\|\boldsymbol{x}_K\times\boldsymbol{p}_{\mu_{\mathrm{r}}}\|_{K_{n}}}\le\frac{1}{n}
    \ \forall n\in\mathbb{N}.
  \end{equation}

  Since \(\mathscr{M}(S_{N_{\mathrm{v}}})\) is bounded and closed, there exists a subsequence \(\{\mathscr{M}(K_{n_{i}})\}_{i\in \mathbb{N}}\) converging to \(\mathscr{M}(K)\), which corresponds to \(K\) as the limit of \(\{K_{n_{i}}\}_{i\in \mathbb{N}}\) in the sense that all vertex coordinates converge. Since:
  \begin{equation}
    \|\nabla \times \boldsymbol{q}_{n}\|_{\tilde{B}}\le\|\nabla \times \boldsymbol{q}_{n}\|_{K_{n}}=1
  \end{equation}

  is bounded, there exists another subsequence of \(\{\nabla \times \boldsymbol{q}_{n_{i}}\}_{i\in \mathbb{N}}\) converging to \(\nabla \times \boldsymbol{q}\).

  By standard polynomial properties, we have:
  \begin{equation}
    1=\|\nabla \times \boldsymbol{q}\|_{K}\le\|\nabla \times \boldsymbol{q}\|_{B}\lesssim\|\nabla \times \boldsymbol{q}\|_{\tilde{B}},
  \end{equation}
  which implies \(\nabla \times \boldsymbol{q}\neq 0\) in \(\tilde{B}\). But, following the same argument as in the proof to \cref{prop:sf}, the fact that last two terms in \cref{eq:pfcontrasf} diminish, implies \(\nabla \times \boldsymbol{q} =  0\) in \(K\) hence in \(\tilde{B}\). This is a contradiction, and thus the claim is proved.
\end{proof}

\end{document}
